\section{Conclusion}

When a Spark predicate reports that two schemas are equal, what it has agreed to about optionality is one of two things and never a third. Over a drift taxonomy generated from \texttt{StructType}'s grammar, the ten schema-equality configurations a stated discovery rule admits over Spark 3.5.6 -- and, by source inspection, through \texttt{v4.2.0} -- either accept a change to any of the three carriers or reject a change to any of them. None discriminates, so no caller can be strict about one carrier and lenient about another. That is the result, and it is a consequence of modelling a three-dimensional property with one \texttt{Boolean} parameter rather than an oversight in any single predicate.

The result matters because the demand comes from outside Spark and the property it is about is one real schemas use. Avro's field-matching rule and a directional containment requirement on the field carrier jointly require what no shipped predicate provides, and a quarter of the Avro schemas in a twenty-repository corpus declare optionality unevenly across the three carriers rather than uniformly.

The gap cannot be closed where it is observed. \texttt{StructField.nullable} carries three producer-side meanings in two values, and Spark's readers collapse two of them, so the claim survives only upstream, in the Scala types where \texttt{Option[A]} and \texttt{A} differ. The artifact presented here follows that boundary rather than arguing with it: the field carrier is compared at compile time and deliberately dropped at runtime through an explicitly named method. Its runtime half is on the failing side of its own finding, which is how the finding was located, and what it says about the artifact is reported alongside what it says about Spark.

It is worth being plain about how much of the empty cell the artifact occupies, because a reader is entitled to ask. One half of it is occupied by a working mechanism: a compile-time check that satisfies both externally sourced requirements, which no shipped predicate does. The other half is occupied by an argument that nothing can occupy it -- an elimination over the candidate runtime rules, each failing for its own stated reason. A mechanism plus an impossibility argument is a smaller contribution than two mechanisms, and it is the contribution the measurement supports. The alternative would have been to ship a runtime check that cannot fire.

What remains unmeasured is whether any of this changes an outcome a practitioner would notice. The characterisation is complete for one property over an enumerated predicate population; the effectiveness claim is not attempted.
